% Generated draft: requires mathematical and source review.
\documentclass[11pt]{article}
\usepackage[margin=25mm]{geometry}
\usepackage{fontspec}
\setmainfont{Latin Modern Roman}
\usepackage{amsmath,amssymb,mathtools,unicode-math}
\setmathfont{Latin Modern Math}
\usepackage{enumitem,needspace,xurl,hyperref}
\setlength{\parindent}{0pt}
\setlength{\parskip}{5pt}
\setlength{\emergencystretch}{3em}
\newcommand{\N}{\mathbb{N}}
\newcommand{\Z}{\mathbb{Z}}
\newcommand{\Q}{\mathbb{Q}}
\newcommand{\R}{\mathbb{R}}
\newcommand{\C}{\mathbb{C}}
\newcommand{\F}{\mathbb{F}}
\newcommand{\A}{\mathbb{A}}
\newcommand{\PP}{\mathbb{P}}
\newcommand{\E}{\mathbb{E}}
\newcommand{\eps}{\varepsilon}
\newcommand{\dd}{\,\mathrm{d}}
\begin{document}
\section*{UnsolvedMath ID: 30004250}
\textbf{Source problem code:} catalogue entry
\subsection{Definitions and mathematical statement}
The source poses a research programme rather than a conjecture with a specified conclusion. Its ground field, characteristic, dimension, deformation bases, and precise convention for \(\mathbb Q\)-Gorenstein families are not specified. The definitions below give a conventional framework for interpreting the question; these choices require confirmation against the source.

Let \(k\) denote the unspecified ground field. A toric variety \(X\) over \(k\) is a normal variety containing an algebraic torus as a dense open subset, with the torus action extending to \(X\). Under the convention used here, a toric Fano variety is projective, its canonical divisor \(K_X\) is \(\mathbb Q\)-Cartier, and \(-K_X\) is ample. A divisor is \(\mathbb Q\)-Cartier if some positive integral multiple is Cartier. The phrase ``with non-Gorenstein singularities'' is interpreted as requiring at least one point \(x\in X\) for which the local ring \(\mathcal O_{X,x}\) is not Gorenstein. Whether the source intends any narrower singularity class is unspecified.

A deformation of \(X\) over a pointed \(k\)-scheme \((S,s)\), with residue field \(k\) at \(s\), consists of a flat proper morphism
\[
\pi\colon\mathcal X\longrightarrow S,
\qquad \mathcal X_s\cong X.
\]
For a coherent sheaf \(\mathcal F\), write \(\mathcal F^{**}=\mathcal H om(\mathcal H om(\mathcal F,\mathcal O),\mathcal O)\). Whenever a suitable relative dualizing sheaf exists, write
\[
\omega_{\mathcal X/S}^{[m]}=
\bigl(\omega_{\mathcal X/S}^{\otimes m}\bigr)^{**}
\qquad(m\geq 1).
\]
One common strong convention for a \(\mathbb Q\)-Gorenstein deformation requires these reflexive powers to be flat over \(S\) and compatible with base change, and requires \(\omega_{\mathcal X/S}^{[r]}\) to be invertible for some positive integer \(r\), locally on the family. On normal spaces where relative canonical divisors are defined, the latter condition expresses that \(K_{\mathcal X/S}\) is \(\mathbb Q\)-Cartier. The source does not specify the hypotheses on families needed for this formulation or whether it intends this base-change convention or a weaker relative \(\mathbb Q\)-Cartier condition.

\textbf{Problem.} Develop a combinatorial and deformation-theoretic description of the \(\mathbb Q\)-Gorenstein deformations of toric Fano varieties having non-Gorenstein singularities. In a setting where toric varieties are described by fans in lattices, the requested combinatorial description should relate the toric data of \(X\) to its \(\mathbb Q\)-Gorenstein deformation theory. The original question is: ``What can be said about Q-Gorenstein deformations of toric Fano varieties whose singularities are not Gorenstein?'' Neither a specific classification target nor a necessary-and-sufficient criterion is supplied. The statement does not impose torus equivariance on the deformations; whether equivariant deformations or general deformations are intended remains unresolved.
\subsection{Short English statement}
Describe, using toric combinatorics and deformation theory, the \(\mathbb Q\)-Gorenstein deformations of toric Fano varieties with non-Gorenstein singularities. The source leaves the ground field, allowed families, precise \(\mathbb Q\)-Gorenstein convention, and expected form of the description unspecified. This is a statement-only research-problem draft requiring mathematical review.
\subsection{Sources}
\begin{itemize}
\item Local UnsolvedMath registry, supplied record with permanent ID \(30004250\), title ``Deformations of Non-Gorenstein Toric Fano Varieties,'' source problem code \texttt{catalogue entry}, in the Oberwolfach Reports Open Problems collection. The record gives proposed year \(2019\) and names no proposer.
\item Supplied primary report citation: \emph{Toric Geometry (2020)}, Oberwolfach Reports, pp.\ 2671--2726, DOI \texttt{10.4171/owr/2019/43}. \url{https://doi.org/10.4171/owr/2019/43}. The supplied background describes Oberwolfach Reports as documenting workshops and their open-problem discussions. The citation's \(2020\) title date and the record's \(2019\) proposed year are retained as supplied.
\item Supplied literature source, labelled ``Toric deformation spaces'': \url{https://doi.org/10.1090/tran/8636}.
\item Supplied literature source, labelled ``Toric Q-Gorenstein singularities'': \url{https://arxiv.org/abs/alg-geom/9403003}.
\item Literature assessment checked at \texttt{2026-08-22}: the supplied curation labels the problem \texttt{partially\_solved} and reports that \(\mathbb Q\)-Gorenstein deformation theory and classifications exist in restricted singularity classes, while no general solution for non-Gorenstein toric Fano varieties was located. This is supplied dated curation derived from the OWR record, not an independently verified literature assessment. Its time-sensitive status and cited claims require independent verification.
\item The registry labels the statement \texttt{corrected\_verified} and reports that mathematical content was preserved while notation, typography, and sentence structure were repaired for publication. These are supplied editorial qualifications, not mathematical verification of this draft.
\end{itemize}
\end{document}
